\section*{Chapitre \ref{ch:m1:exchange-traded-funds} --- Fonds indiciels cotés}

\begin{solution}{exo:m1:exchange-traded-funds:1}
$82.31/82.40 - 1 = -10.9$ points de base~: une décote, hors de la bande de 9.
Le \omterm{def:m1:exchange-traded-funds:ap}{participant autorisé} achète des parts d'ETF en bourse, vend le panier
(à découvert si besoin), et rachète les parts de l'ETF auprès du fonds contre le panier,
ce qui clôt sa position à découvert. Il gagne environ 1.9 points de base après coûts.
\end{solution}

\begin{solution}{exo:m1:exchange-traded-funds:2}
L'unité vaut $50\,000 \times 40 = 2$ millions de dollars~; $500/2\,000\,000 = 2.5$
points de base. Sur dix unités, elle vaut 0.25 point de base~: les frais étant fixés
par ordre, la bande est donc plus étroite pour ceux qui créent en volume, et les petits fonds
ayant de petites créations ont des bandes plus larges.
\end{solution}

\begin{solution}{exo:m1:exchange-traded-funds:3}
Le fonds rapporte $-2 \times -4\% = +8\%$~; ses actifs nets passent à 864~millions de dollars.
$\Delta = \beta(\beta-1)Nr = 6 \times 800 \times (-0.04) = -192$ millions de dollars~:
l'exposition doit diminuer de 192~millions de dollars. Pour un fonds inverse, «~moins d'exposition~»
signifie une position à découvert plus importante~: il vend 192~millions de dollars de l'indice, après une baisse. Vérification~:
la position à découvert valait 1\,600, elle a diminué avec l'indice pour atteindre 1\,536, et la cible est $2
\times 864 = 1\,728$.
\end{solution}

\begin{solution}{exo:m1:exchange-traded-funds:4}
$2\times$~: $1.20 \times 0.8182 - 1 = -1.82\%$. $3\times$~: $1.30 \times 0.7273
- 1 = -5.45\%$. $-1\times$~: $0.90 \times 1.0909 - 1 = -1.82\%$. Le
fonds à trois fois l'indice perd le plus~: la perte croît comme $\beta(\beta-1)$, soit
6 pour $\beta = 3$ et 2 pour $\beta = 2$ comme pour $\beta = -1$.
\end{solution}

\begin{solution}{exo:m1:exchange-traded-funds:5}
$\tfrac12\beta(\beta-1)\sigma^2 = 0.16$, $0.48$ et $0.48$~: facteurs
$\mathrm{e}^{-0.16} = 0.852$ pour $\beta = 2$ et $\mathrm{e}^{-0.48} = 0.619$
pour $\beta = 3$ comme pour $\beta = -2$. Une perte de 10~\%~: $0.48\,T = -\ln 0.9 =
0.105$, donc $T = 0.22$ année, soit environ 55 jours de bourse.
\end{solution}

\begin{solution}{exo:m1:exchange-traded-funds:6}
L'écart entre la valeur et la \omterm{def:m1:exchange-traded-funds:nav}{valeur liquidative} vaut $6 \times 0.8^{k+1}$ le $k$-ième jour après
la chute. \omterm{def:m1:exchange-traded-funds:nav}{Valeur liquidative}~: 98.80, 97.84, 97.07, 96.46. Décote~: 4.86~\%, 3.92~\%, 3.16~\%,
2.55~\%. Elle passe sous 0.5~\% le onzième jour après le jour de la chute.
Un observateur patient voit une décote de plusieurs pour cent «~se refermer~» pendant deux
semaines alors que le cours de l'ETF ne bouge pas du tout.
\end{solution}

\begin{solution}{exo:m1:exchange-traded-funds:7}
Avec une bande de 7 points de base, un \omterm{def:m1:exchange-traded-funds:ap}{participant autorisé} agit à 1\,638 des
5\,000 pas, et la \omterm{def:m1:exchange-traded-funds:nav}{prime} absolue moyenne vaut 3.28 points de base. Avec $c_B =
2$, la bande vaut 5~: 2\,219 actions et une \omterm{def:m1:exchange-traded-funds:nav}{prime} absolue moyenne de 2.42. L'investisseur
qui achète ou vend à un moment quelconque paie, en moyenne, un écart plus faible
par rapport à la fair value~; des marchés sous-jacents moins chers rendent les ETF plus serrés, et
une activité plus soutenue sur le marché primaire est le signe d'un produit plus serré, et non
d'un produit sous tension.
\end{solution}

\begin{solution}{exo:m1:exchange-traded-funds:8}
La «~décote~» compare un cours en temps réel à une \omterm{def:m1:exchange-traded-funds:nav}{valeur liquidative} calculée à partir d'obligations
évaluées de façon obsolète (voir la \cref{prop:m1:exchange-traded-funds:stale}). Elle peut
se résorber entièrement par la baisse de la \omterm{def:m1:exchange-traded-funds:nav}{valeur liquidative} jusqu'au cours de l'ETF, auquel cas l'acheteur de l'ETF
ne gagne rien, ou perd si le marché continue de baisser. Pour capter une décote réelle,
il faut acheter l'ETF \emph{et vendre le panier} aux cours utilisés
dans la \omterm{def:m1:exchange-traded-funds:nav}{valeur liquidative}, puis racheter~; sur les obligations à haut rendement dans un marché baissier,
personne n'achètera le panier à ces cours, c'est précisément pour cela que la décote existe.
«~Elle se résorbe toujours~» est vrai et sans intérêt~; «~sans risque~» est faux.
\end{solution}

\begin{solution}{pb:m1:exchange-traded-funds:1}
\textbf{1.} $+3\%$ et $-3\%$~: 8.24~milliards et 2.91~milliards de dollars.
\textbf{2.} $6 \times 8 \times 0.01 = 0.48$ milliard de dollars et $12 \times 3 \times
0.01 = 0.36$ milliard de dollars, tous deux à l'achat~: 0.84~milliard de dollars.
\textbf{3.} Les dix dernières minutes échangent 30~milliards de dollars~: 2.8~\%.
\textbf{4.} $\sigma_{10} = 1\% \times \sqrt{10/390} = 0.16\%$~; impact $0.7
\times 0.16\% \times \sqrt{0.028} = 1.9$ points de base. Imperceptible.
\textbf{5.} $-15\%$ et $+15\%$~: 6.8~milliards et 3.45~milliards de dollars.
\textbf{6.} $-2.4$ milliards de dollars et $-1.8$ milliards de dollars~: 4.2~milliards de dollars à vendre.
\textbf{7.} 14~\%.
\textbf{8.} $\sigma_{10} = 0.48\%$~; $0.7 \times 0.48\% \times \sqrt{0.14} =
12.6$ points de base.
\textbf{9.} La transaction vaut $(6 \times 8 + 12 \times 3) = 84$ milliards par unité de
rendement~; un supplément de $0.126\%$ ajoute 0.11~milliard de dollars, soit 2.5~\% de plus. La rétroaction
existe et converge aussitôt à cette taille~; elle ne convergerait pas si les fonds étaient dix
fois plus grands par rapport à la clôture.
\textbf{10.} $0.95 \times 1.04 \times 0.94 \times 1.05 \times 1.025 - 1 =
-0.05\%$~: quasiment nul.
\textbf{11.} $0.85 \times 1.12 \times 0.82 \times 1.15 \times 1.075 - 1 =
-3.49\%$, contre $-0.14\%$.
\textbf{12.} $1.15 \times 0.88 \times 1.18 \times 0.85 \times 0.925 - 1 =
-6.11\%$, contre $+0.14\%$.
\textbf{13.} $-4.80\%$ sur une semaine où l'indice n'a pas bougé.
\textbf{14.} Les deux fonds subissent une érosion au rythme $\tfrac12\beta(\beta-1)\sigma^2$~: 3
et 6 fois la variance de l'indice. Les expositions à l'indice se compensent~; les érosions s'additionnent.
La paire constitue une position courte sur la variance réalisée, payée par cinq jours de
mouvements de 4--6~\%. (La transaction symétrique, vendre les deux fonds, n'est longue de rien et
courte de l'érosion~: elle capte la variance mais est exposée à une tendance, dans laquelle une jambe
se capitalise sans limite, et au coût d'emprunt des parts.)
\textbf{15.} Les \omterm{def:m1:what-a-trading-firm-does:market-maker}{teneurs de marché} et les \omterm{def:m1:what-a-trading-firm-does:prop}{sociétés de trading pour compte propre}, qui connaissent les actifs des fonds
(publiés chaque jour) et le rendement de la journée (public), donc la taille et le sens de
la transaction, bien avant la clôture. Ils achètent ou vendent avant elle et fournissent
la contrepartie dans l'enchère de clôture (au \cref{ch:m1:auctions})~; leur
concurrence maintient l'impact près de l'estimation de la question 8.
\textbf{16.} $\beta(\beta-1)$ vaut 12 contre 6. Après une hausse, les actifs du fonds inverse
diminuent tandis que sa position à découvert gagne en valeur~: ces deux effets font monter
son levier, alors que pour le fonds long ils se compensent en partie.
\textbf{17.} $-33.3\%$ en une journée. Les coupe-circuits de marché américains suspendent
la négociation pour la journée lorsque le S\&P 500 a chuté de 20~\%, ce qui limite la perte
d'un jour d'un fonds sur un indice américain large en dessous de ce niveau~; les fonds sur une action isolée,
un secteur ou la volatilité n'ont aucune protection de ce type, et leurs prospectus avertissent
qu'un investisseur peut perdre l'intégralité de son investissement en une seule journée.
\textbf{18.} La transaction de rééquilibrage, rapportée au volume de clôture~: 14~\% ici,
2.8~\% le jour calme. Elle croît comme les actifs, multipliés par $\beta(\beta-1)$ et par le
mouvement, elle est donc maximale précisément les jours où la clôture
est déjà tendue.
\textbf{19.} \textbf{4.2~milliards de dollars à vendre.}
\textbf{20.} Pour quelqu'un qui veut un levier pour un jour ou quelques jours, sans compte
de marge, et qui sait qu'au-delà de cet horizon le produit est une
position sur la trajectoire, et pas seulement sur la destination.
\end{solution}

\begin{solution}{iq:m1:exchange-traded-funds:1}
Les \omterm{def:m1:exchange-traded-funds:ap}{participants autorisés} peuvent échanger les parts de l'ETF contre le panier sous-jacent
auprès du fonds, dans les deux sens, à la \omterm{def:m1:exchange-traded-funds:nav}{valeur liquidative}. Si l'ETF est cher, ils le vendent
achètent le panier et créent des parts~; s'il est bon marché, c'est l'inverse. Le cours reste
à l'intérieur d'une bande fixée par le coût de négociation du panier, de l'ETF, des frais et du financement.
La plupart du temps, ce sont les \omterm{def:m1:what-a-trading-firm-does:market-maker}{teneurs de marché} qui couvrent leur position avec des contrats à terme ou avec le panier,
et les créations ne constituent que le résidu net.
\iqlookfor{les deux sens, la bande plutôt que «~égal à la \omterm{def:m1:exchange-traded-funds:nav}{valeur liquidative}~», et
le fait que le fonds lui-même ne négocie pas.}
\end{solution}

\begin{solution}{iq:m1:exchange-traded-funds:2}
Presque jamais exactement. Le fonds livre deux fois chaque \emph{rendement quotidien}~; sur un an, le résultat vaut le carré du rapport
des indices, multiplié par $\exp(-\sigma^2 T)$. Avec une volatilité de 20~\%~: $1.21 \times 0.961 = 1.163$, soit
une hausse d'environ 16~\% avant frais et financement~; dans une année
de tendance régulière, elle peut dépasser
20~\%.
\iqlookfor{le mot «~quotidien~», un chiffre, et la remarque que les tendances aident.}
\end{solution}

\begin{solution}{iq:m1:exchange-traded-funds:3}
Non. La \omterm{def:m1:exchange-traded-funds:nav}{valeur liquidative} utilise la clôture de Tokyo, vieille de dix heures. Depuis, les contrats à terme sur l'indice
japonais, le yen et les marchés américains ont bougé~; le cours de l'ETF est l'estimation par le marché
de l'ouverture de Tokyo. Il faut comparer l'ETF à une fair value construite à partir d'instruments ouverts maintenant.
Un arbitrage supposerait de négocier
le panier aux cours périmés, ce qui est impossible.
\iqlookfor{refus du mot «~arbitrage~», et une construction concrète
de fair value.}
\end{solution}

\begin{solution}{iq:m1:exchange-traded-funds:4}
Je suis à découvert de 50~millions de dollars sur l'ETF. J'achète immédiatement
des contrats à terme sur indice pour le même delta~: une transaction peu coûteuse, qui laisse
un risque de base entre le contrat et le panier. Au cours de la journée, d'autres clients me vendent
des parts et réduisent la position. Près de la clôture, pour ce qui reste, j'achète soit
le panier (un program trade, souvent dans l'enchère de clôture, afin que mon coût corresponde à la \omterm{def:m1:exchange-traded-funds:nav}{valeur liquidative}),
et je crée des parts pour couvrir le découvert, en vendant les contrats à terme dans le même temps, soit je garde
la position courte sur l'ETF contre contrats à terme pour la nuit, si emprunter l'ETF est bon marché et que la création
ne vaut pas les frais. Mon avantage est la fourchette facturée, diminuée des coûts de couverture,
des frais et du financement.
\iqlookfor{couvrir d'abord, créer ensuite~; conscience du strike de la \omterm{def:m1:exchange-traded-funds:nav}{valeur liquidative}
à la clôture et de la compensation des positions d'inventaire.}
\end{solution}

\begin{solution}{iq:m1:exchange-traded-funds:5}
Actifs $N$, exposition $\beta N$. Après un rendement $r$~: actifs $N(1+\beta r)$,
exposition $\beta N(1+r)$, cible $\beta N(1+\beta r)$~; la transaction vaut $\beta(\beta-1)Nr$.
Le coefficient est positif aussi pour un $\beta$ négatif~: les fonds inverses achètent après
les hausses et vendent après les baisses, comme les \omterm{def:m1:exchange-traded-funds:leveraged}{ETF à effet de levier}, et un fonds $-1\times$
négocie autant qu'un fonds $2\times$.
\iqlookfor{une dérivation propre en trois lignes et le signe du résultat pour les
fonds inverses.}
\end{solution}

\begin{solution}{iq:m1:exchange-traded-funds:6}
L'ETF est mal évalué~: les bilans des \omterm{def:m1:exchange-traded-funds:ap}{participants autorisés} sont contraints, si bien que
la bande s'est élargie~; les vendeurs forcés utilisent l'ETF parce que c'est le seul
actif liquide. La \omterm{def:m1:exchange-traded-funds:nav}{valeur liquidative} est mal évaluée~: les obligations n'ont pas été échangées, les valorisations sont en retard, et
l'ETF prend l'avance. Éléments de preuve~: les cours des obligations échangées, comparés à leurs valorisations~;
le comportement de la \omterm{def:m1:exchange-traded-funds:nav}{valeur liquidative} dans les jours suivants (descend-elle vers
l'ETF~?)~; les indices de credit default swaps et les contrats à terme comme fair values indépendantes~;
l'activité de rachat (si les \omterm{def:m1:exchange-traded-funds:ap}{participants autorisés} rachètent massivement, la décote est réelle pour eux)~; et
si les paniers de rachat se négocient à des prix inférieurs à la
valorisation.
\iqlookfor{les deux cas argumentés, et des données qui discriminent, en particulier la
convergence ultérieure de la \omterm{def:m1:exchange-traded-funds:nav}{valeur liquidative} vers l'ETF.}
\end{solution}
